\section{DOM: Traversing, Membuat Node, dan Menghapus Elemen}

\textbf{DOM traversing} memungkinkan navigasi dari satu elemen ke elemen lain melalui hubungan parent-child-sibling. Properti traversing: \code{parentElement} (naik ke parent), \code{children} (turun ke anak-anak elemen), \code{firstElementChild}, \code{lastElementChild}, \code{nextElementSibling}, \code{previousElementSibling}. Versi tanpa ``Element'' (\code{childNodes}, \code{firstChild}) juga menyertakan text nodes dan komentar \cite{jsInfo}.

\begin{table}[htbp]
\centering
\small
\begin{tabular}{|P{4.2cm}|P{3.2cm}|P{2.6cm}|}
\hline
\textbf{Operasi} & \textbf{Metode} & \textbf{Catatan} \\
\hline
Buat elemen & \raggedright\code{document.}\linebreak[0]\code{createElement(tag)} & Belum ada di DOM \\
\hline
Isi teks & \code{el.textContent = "..."} & Aman dari XSS \\
\hline
Tambah anak & \raggedright\code{parent.}\linebreak[0]\code{appendChild(el)} & Di akhir parent \\
\hline
Sisip sebelum & \raggedright\code{parent.}\linebreak[0]\code{insertBefore(el, ref)} & Sebelum elemen ref \\
\hline
Hapus elemen & \code{el.remove()} & Hapus dari DOM \\
\hline
Clone & \code{el.cloneNode(true)} & deep clone \\
\hline
\end{tabular}
\caption{Operasi DOM untuk membuat dan menghapus elemen}
\end{table}

\begin{catatan}
\textbf{Minimalisir Manipulasi DOM.} Setiap perubahan DOM memicu layout recalculation oleh browser (reflow). Untuk menambah banyak elemen, kumpulkan di \code{DocumentFragment} terlebih dahulu lalu tambahkan sekali: \code{const frag = document.createDocumentFragment(); frag.appendChild(el1); frag.appendChild(el2); parent.appendChild(frag);}. Ini jauh lebih efisien daripada \code{appendChild()} berulang \cite{mdnJsGuide}.
\end{catatan}

